% Candidate for v0.14; not a canonical edit or publication approval.
% RESULT.json SHA256 7c9ff31f20b4ce7ee65e146a13ba1ccdda18ae04c07757e7ad1abf56b4702e88
% RESULT-REVIEW.txt SHA256 4bb05fb4bcaa8c187c7fd850c0d2908283bedf98b3e1d494e5e3676658373e33
\paragraph{Changing the complete pulse shape.}
A subsequent fixed comparison used the same nominal linear radial
$\ell=1$ evolution, $R=20$, $T=6$ and $h=0.04,0.02$.
It compared the annular drive above with a mode-adapted drive, using
single pulses and equal-sign pairs separated by $\pi/\rho$.
The new drive retains both sideband terms of a quintic turn-on construction
and has a spatial cutoff equal to one for $r\leq4$ and zero for $r\geq8$.
The reference instead uses the annulus $1<r<2$ and a $\sin^2$ temporal
window with one rotating frequency. Thus both spatial and temporal
waveforms change; the comparison does not isolate a spatial cause.

Each single pulse was normalized to the same reduced-source norm
\begin{equation}
 B_{\mathrm{src}}=\int\!\int |s(r,t)|^2\,\dd r\,\dd t,
 \qquad B_{\mathrm{src}}\simeq1.35281\times10^{-3}
 \quad(h=0.02).
\end{equation}
Pairs have twice this budget. The integrated budgets agreed with their
prescribed values within $6\times10^{-9}$ relatively on both grids.
This norm is neither mechanical work nor injected energy or charge,
and equal norms do not imply equal actuator support or hardware cost.
For the single pulse, the fine-grid diagnostic $|d(T)|^2/B_{\mathrm{src}}$
is about $32$ times larger for the shaped drive
(Table~\ref{tab:source-shaping}); this is not an energy-efficiency ratio
or an optimality result.

\begin{table}[htbp]\centering\small
\caption{Fixed source-shaping comparison at $T=6$, $h=0.02$.
Here $p=\partial_tq$, and ``remainder'' means subtraction of the nominal
mode reconstruction. Field norms are $L^2(\dd r)$ norms, not energies;
the pair rows give the total remaining field and velocity norms.}
\label{tab:source-shaping}
\begin{tabular}{lrr}\toprule
Diagnostic&Annular drive&Shaped drive\\\midrule
Single-pulse $|d|^2/B_{\mathrm{src}}$&$0.001214$&$0.03872$\\
Single-pulse $q$ remainder norm&$0.007992$&$0.002596$\\
Single-pulse $p$ remainder norm&$0.01502$&$0.004954$\\
Pair $\|q\|$&$0.01193$&$0.001055$\\
Pair $\|p\|$&$0.02153$&$0.001930$\\\bottomrule
\end{tabular}
\end{table}

The shaped pair leaves approximately $8.8\%$ and $9.0\%$ of the
reference pair's field and velocity norms, respectively; both remain
nonzero. Nominal mode subtraction is not an orthogonal energy
partition. Both grids passed the fixed comparison and defect-balance
checks, but the global second-component nominal eigenmode residual norm increases
from approximately $0.529$ to $1.495$ on refinement. The successful
paired-defect balance therefore does not establish global eigenfunction
convergence. These finite-box, short-time results show reduced residual
motion for the specified composite drive; they establish neither full
stillness of the field nor removal of the charged background, nonlinear
stability or a lifetime. This original two-source comparison did not
separate changes of spatial and temporal design. The subsequently
specified cross-comparison below addresses that limitation within a
fixed actuator family.

% Reviewed stored-profile attribution, incorporated in draft 0.17.
% RESULT 03e1c3129984f7b96713f901e7e7767b5f32d4bb6b33afa380178d34dc8ffbb7
% RESULT-REVIEW 3b665e564d1471fcfbf80cb42d118615c011bfb7b4a2cffea75522ca27619b9c
A subsequent stored-profile audit identifies the origin of the growing
second-component defect in this box. The exported reference mode has
$v(20)=0.004228503237$, whereas the evolution operator imposes a zero
Dirichlet value. For the three-point stencil this discrepancy adds
$b_j=\delta_{j,n-1}v(20)/h^2$ to the continued-ghost residual, giving
$\|b\|_h^2=|v(20)|^2/h^3$. Its squared norm is $0.2793787442$ and
$2.2350299533$ on the two grids; the continued residual squared norms
are $1.52\times10^{-8}$ and $9.50\times10^{-10}$, with signed cross
terms $-9.74\times10^{-10}$ and $-4.84\times10^{-10}$.
The boundary discrepancy therefore dominates this component's growing
residual. It does not dominate the first component. This is an algebraic
attribution, not a modification of the evolution boundary condition or a
new eigenfunction calculation. The original global defect remains valid
for the original operator. Two decreasing interior residuals neither
certify the exported mode nor bound the integrated influence of the
boundary discrepancy on the pulse response.

\begin{figure}[htbp]\centering
\includegraphics[width=.98\linewidth]{source-shaping.pdf}
\caption{Endpoint diagnostics from the fixed spatiotemporal pulse comparison.
Both spatial support and temporal waveform change. The common source norm
is not energy or work. The shaped source increases the selected projection
per source norm and reduces the residual response in this comparison;
the complete field does not vanish. Bars show the stored fine-grid result,
not a new evolution or a certified continuum limit.}
\label{fig:source-shaping}
\end{figure}

% Factorial RESULT477978704a19d5b2f2b474c41cf5917367335c521902e0f712303f7a563ec263
% Non-author RESULT-REVIEW66fe64890bbbc3568c19a33610968cd7360f68462ea8f7fb7b467b900dd02a28
\paragraph{A fixed factorial actuator comparison.}
The new comparison uses two spatial pairs, $A=(w,0)$ and
$B=(\zeta u,\zeta v)$, where $w$ is the preceding annulus profile,
$\zeta$ the cutoff defined above (one below $r=4$ and zero above $r=8$),
and $u,v$ the stored nominal sidebands.
Two temporal pairs are prescribed on $0<t<1$:
\begin{align}
 T_0&=\bigl(\sin^2(\pi t)e^{i\rho t},
                  \sin^2(\pi t)e^{-i\rho t}\bigr),\\
 T_1&=\bigl([g''+2i(\rho+\omega)g']e^{i\rho t},
                  [g''-2i(\rho-\omega)g']e^{-i\rho t}\bigr),
 \qquad g=10t^3-15t^4+6t^5.
\end{align}
Both temporal pairs vanish outside $0<t<1$; $g$ is extended constantly
as zero before the pulse and one afterward.
Each complex source is the component sum
$s_{ij}=\mathcal N_{ij}(R_{i1}T_{j1}+R_{i2}T_{j2})$,
where $R_i=A$ or $B$. There is no conjugation in this sum; the second
physical equation receives $s_{ij}^*$.
The positive factor $\mathcal N_{ij}$ fixes the same single-pulse
$B_{\mathrm{src}}$ \emph{after} summation, retaining component interference.
The two old corners $A/T_0$ and $B/T_1$ are implementation controls;
$A/T_1$ and $B/T_0$ are the new crossed drives. Each uses both a single
pulse and an equal-sign pair at the unchanged delay $\pi/\rho$.
The nominal operator, $R=20$, $T=6$, $h=0.04,0.02$ and time step $h/8$
remain fixed. Both grids passed the predeclared numerical checks.

\begin{table}[htbp]\centering\small
\caption{Four fixed normalized actuator designs at $T=6$, $h=0.02$.
The projection diagnostic is for a single pulse; the final two columns
are total remaining pair norms, not mode-subtracted norms or energies.}
\label{tab:source-factorial}
\begin{tabular}{lrrr}\toprule
Design&$|d_{\mathrm{single}}|^2/B_{\mathrm{src}}$&Pair $\|q\|$&Pair $\|p\|$\\\midrule
$A/T_0$&$0.001214$&$0.01193$&$0.02153$\\
$A/T_1$&$0.0004100$&$0.006390$&$0.01426$\\
$B/T_0$&$0.1586$&$0.02024$&$0.005692$\\
$B/T_1$&$0.03872$&$0.001055$&$0.001930$\\\bottomrule
\end{tabular}
\end{table}

Within this family, $B/T_0$ gives the largest single-pulse projection
diagnostic, whereas $B/T_1$ leaves the smallest pair norms in both $q$
and $p$. Strong projected excitation and small residual field are thus
different objectives. The spatial-pair change $A\to B$ increases the
pair $q$ norm at $T_0$ but decreases it at $T_1$: its effect depends on
the temporal partner. These are conditional design comparisons, not
universal additive spatial and temporal effects. In particular $A$
omits a spatial component, $T_1$ changes sideband weights and phase
structure, and normalization depends on both factors. Equal source norm
still implies neither equal work nor energy efficiency. All four pairs
suppress their selected projection under the fixed criterion, while their
full responses differ and remain nonzero. The nominal-mode defect and
finite-box limitations above remain. That fixed-phase comparison did not
test phase-error robustness; it established neither exact preparation,
complete stillness nor a lifetime.

\begin{figure}[htbp]\centering
\includegraphics[width=.98\linewidth]{factorial-comparison.pdf}
\caption{Stored endpoint diagnostics for the fixed $2\times2$ actuator
family. The largest single-pulse projection per source norm occurs at
$B/T_0$, while $B/T_1$ gives the smallest remaining pair field and
velocity norms. The source is normalized after adding its two components;
these comparisons are not energy efficiencies or a universal separation
of independent spatial and temporal mechanisms. No additional evolution
or phase-error test is represented by this figure.}
\label{fig:source-factorial}
\end{figure}

% Phase RESULTf492a34c52a608a36d3d6054c1342cc6bf3dd963021e3576d5d6447423a25301
% Reviewed ea5b091239d228e8da5ac68c60108c238c3052a2ee4272f9012366e61459ec37
% Endpoint SUMMARY3329b5213520fe157b71c0a79d530a36e03ce9a644fd13d48cc71031e0c53aae
\paragraph{A constant relative pulse phase.}
A further bounded test changes only the additional phase $\theta$ of the
second pulse, keeping the delay and the carrier-phase correction fixed.
The coupling to $q^*$ makes the nominal Bogoliubov evolution real-linear,
not complex-linear: the response to $i s$ need not equal $i q$.
For each of the four fixed designs, one new delayed-$i s$ response
therefore completes the real basis. At the same endpoint,
\begin{equation}
 q_\theta=q_{\mathrm{single}}
   +\cos\theta\,(q_{\mathrm{pair},0}-q_{\mathrm{single}})
   +\sin\theta\,q_{\mathrm{delayed},90}.
 \label{eq:phase-real-basis}
\end{equation}
The same identity holds for $p$ and the real-linear projection $d$.
Four directly evolved $\theta=\pi/4$ pairs check this reconstruction:
the largest reported relative discrepancy in $q,p,d$ across both grids
was $1.14\times10^{-13}$. This tests a known linear identity, not a new
physical effect. Other constant phases below are evaluations of stored
endfields, not additional time evolutions. The setting remains
$R=20$, $T=6$, $h=0.04,0.02$.

\begin{table}[htbp]\centering\small
\caption{Selected constant-phase endpoints for $B/T_1$ at $T=6$,
$h=0.02$. Norms are total remaining radial fields, not energy.
The phase values were fixed before the endpoint analysis.}
\label{tab:phase-selected}
\begin{tabular}{rrr}\toprule
$\theta$&$\|q\|$&$\|p\|$\\\midrule
$0^\circ$&$0.001055$&$0.001930$\\
$-5^\circ$&$0.001062$&$0.002068$\\
$+5^\circ$&$0.001423$&$0.002462$\\
$-15^\circ$&$0.001986$&$0.003720$\\
$+15^\circ$&$0.002562$&$0.004413$\\
$45^\circ$&$0.006301$&$0.01125$\\\bottomrule
\end{tabular}
\end{table}
These values need not be symmetric about zero phase. Cancellation of the
selected coefficient and minimization of the full field norm are different
conditions; no new minimum search was performed.

For a uniformly distributed phase that is \emph{constant within each
pulse pair}, write any endpoint field as $a+b\cos\theta+c\sin\theta$.
Its root mean square norm is exactly
\begin{equation}
 \sqrt{\mathbb E_\theta\|a+b\cos\theta+c\sin\theta\|^2}
 =\sqrt{\|a\|^2+\tfrac12(\|b\|^2+\|c\|^2)}.
\end{equation}
The Gram-matrix evaluation was also checked against eight equally spaced
phases. Table~\ref{tab:phase-rms} keeps all four designs in the comparison.
\begin{table}[htbp]\centering\small
\caption{Uniform constant-phase RMS endpoint norms, $T=6$, $h=0.02$.
These are square roots of mean squared norms, not mean norms.}
\label{tab:phase-rms}
\begin{tabular}{lrr}\toprule
Design&$\sqrt{\mathbb E_\theta\|q\|^2}$&$\sqrt{\mathbb E_\theta\|p\|^2}$\\\midrule
$A/T_0$&$0.01128$&$0.02168$\\
$A/T_1$&$0.007123$&$0.01447$\\
$B/T_0$&$0.02552$&$0.04088$\\
$B/T_1$&$0.01106$&$0.01954$\\\bottomrule
\end{tabular}
\end{table}
Although $B/T_1$ leaves smaller pair norms than $A/T_1$ at zero phase,
its two RMS norms are larger in this ensemble. Its fixed-phase advantage
is thus conditional on phase preparation.
The comparison fixes source norm, not the previously excited target
amplitude: $A/T_1$ has a much smaller single-pulse projection diagnostic
($0.0004100$ versus $0.03872$ for $B/T_1$). Its lower RMS therefore does
not establish a better storage or erasure protocol at equal preparation.
An RMS ranking is not necessarily
an ordering of mean norms. This ensemble is neither a time average nor a
simulation of fluctuating environmental noise. The constant phase does
not change the disjoint pulses' source budget $2B_{\mathrm{src}}$, but that
budget remains distinct from work or energy. The global nominal eigenmode
residual limitation remains unchanged; neither exact BIC preparation,
complete stillness nor nonlinear or long-time robustness follows.

\begin{figure}[htbp]\centering
\includegraphics[width=.98\linewidth]{phase-response.pdf}
\caption{Constant-phase endpoint reconstruction for all four fixed
actuator designs, from their stored real-linear response bases.
The two grids compare nominal finite-box observables; plotted phase
samples are not independent time evolutions. The additional phase is
constant within a pulse pair, not a time-dependent noise source.
Field norms and uniform-phase RMS values do not measure energy cost.}
\label{fig:phase-response}
\end{figure}

\paragraph{Quadratic endpoint energy and the rotating generator.}
A further analysis uses the stored single- and pair-pulse endpoints,
without repeating the time evolution. With $p=q_t$ and the same
$h$-weighted inner product and Dirichlet operator as the evolution,
define
\begin{align}
 H_2&=\|p\|_h^2+\langle q,Kq\rangle_h
       +\operatorname{Re}\,h\sum_j C_j q_j^2,
       \label{eq:endpoint-generator}\\
 Q_2&=2\operatorname{Im}\langle q,p\rangle_h
       +2\omega\|q\|_h^2,\qquad E_2=H_2+\omega Q_2.
       \label{eq:endpoint-energy}
\end{align}
Here $K=-\Delta_h+2/r^2+1-4f^2+(9/2)f^4-\omega^2$ and
$C=-2f^2+3f^4$. There is no factor $1/2$: $H_2$ is the coefficient
of second order in $E-\omega Q$ for the stated action convention,
using one physical complex field rather than counting its conjugate twice.
This rotating quadratic generator need not be positive.

\begin{table}[htbp]\centering
\caption{Stored endpoint coefficients at $T=6$, $h=0.02$.
All entries are in units of $10^{-4}$, rounded to four significant digits.
Within a fixed pulse count the designs have equal source-L2 budgets;
a pair has twice the single-pulse budget.}
\label{tab:endpoint-energy}
\begin{tabular}{lrrrr}\toprule
Design&$H_2$ single&$H_2$ pair&$E_2$ single&$E_2$ pair\\\midrule
$A/T_0$&4.782&9.315&5.958&11.51\\
$A/T_1$&2.668&5.180&3.614&6.910\\
$B/T_0$&8.917&1.382&6.467&8.608\\
$B/T_1$&2.129&0.05203&1.201&0.03759\\\bottomrule
\end{tabular}
\end{table}
The distinction matters even for the direction of change:
$B/T_0$ decreases $H_2$ while increasing $E_2$ from single to pair.
Its $Q_2$ changes from approximately $-2.821\times10^{-4}$ to
$8.319\times10^{-4}$. For $B/T_1$ both endpoint coefficients decrease;
for the two $A$ designs both increase. These directions also occur on
$h=0.04$, but two-grid differences are not validated error bounds.

Equal source-L2 budgets do not imply equal work, and the single-to-pair
comparison itself is not a same-budget comparison. For the fixed
semidiscrete model, $\dot H_2=2\operatorname{Re}\langle p,S\rangle_h$.
Thus $H_2(T)$ at zero initial data predicts net rotating-generator work
by an identity; that endpoint-only analysis did not measure an independent
time-integrated power balance. The subsequent test below adds this observable.
The quadratic laboratory source power additionally contains
$2\omega\operatorname{Im}\langle q,S\rangle_h$, while $\dot E_2$
also contains the background-exchange term
$2\omega\operatorname{Im}\,h\sum_j C_jq_j^2$.
In particular, $E_2$ is a coefficient on the linear ansatz, not the complete
laboratory energy change including a second-order background response.
No positive or gross actuator work, energy efficiency, or full shutdown
of the Q-ball is established by Table~\ref{tab:endpoint-energy}.
\paragraph{Time-integrated source work in the nominal linear model.}
A subsequent bounded run retains the same eight single/pair arms,
$R=20$, $T=6$, $h=0.04,0.02$ and $\Delta t=h/8$, and adds power
accumulators at the four Runge--Kutta stages. With the conventions of
Eqs.~\eqref{eq:endpoint-generator}--\eqref{eq:endpoint-energy}, these
integrate
\begin{align}
 P_H&=2\operatorname{Re}\langle p,S\rangle_h,&
 J_Q&=2\operatorname{Im}\langle q,S\rangle_h,&
 X_Q&=2\operatorname{Im}\,h\sum_jC_jq_j^2,\\
 P_{E,\mathrm{src}}&=P_H+\omega J_Q,&
 P_{E,\mathrm{ex}}&=\omega X_Q.
\end{align}
The independently evaluated state coefficients satisfy the numerical
balances $\Delta H_2=\int P_H\,\dd t$,
$\Delta Q_2=\int(J_Q+X_Q)\,\dd t$, and
$\Delta E_2=\int(P_{E,\mathrm{src}}+P_{E,\mathrm{ex}})\,\dd t$.
Unlike the preceding endpoint-only analysis, these integrals use the
actual time-dependent fields and source at the integration stages.
The exchange term remains active after the source ends.

The first pulse is identical in the single and pair arms; the first and
second pulse windows are disjoint. Their trajectories therefore agree before
the second pulse. Subtracting their independently integrated total source
works gives the second pulse's net work on the already excited field:
\begin{equation}
 W_{H,2}=W_{H,\mathrm{pair}}-W_{H,\mathrm{single}},\qquad
 W_{E,\mathrm{src},2}=W_{E,\mathrm{src},\mathrm{pair}}
                       -W_{E,\mathrm{src},\mathrm{single}}.
\end{equation}
This is not the work of an isolated second pulse applied to a fresh
zero state, nor merely a relabeling of endpoint energy differences.

\begin{table}[htbp]\centering
\caption{Net second-pulse source work inferred from time-integrated
pair-minus-single works at $h=0.02$. Units are $10^{-4}$, as in
Table~\ref{tab:endpoint-energy}; entries are rounded to four significant
digits. Negative values denote net extraction in the specified quadratic
source-work convention, not efficiency of a physical actuator.}
\label{tab:second-pulse-work}
\begin{tabular}{lrr}\toprule
Design&$W_{H,2}$&$W_{E,\mathrm{src},2}$\\\midrule
$A/T_0$&$4.533$&$4.207$\\
$A/T_1$&$2.512$&$3.602$\\
$B/T_0$&$-7.535$&$5.939$\\
$B/T_1$&$-2.077$&$-1.122$\\\bottomrule
\end{tabular}
\end{table}
Thus the second $B/T_1$ pulse is a net source-work sink in both
conventions. In contrast, $B/T_0$ removes rotating-generator work while
its laboratory source-work coefficient is positive. The change in $E_2$
also includes the difference of integrated exchange terms; it is not
identical to $W_{E,\mathrm{src},2}$.
Positive and negative parts of the total spatially integrated power
were separately accumulated, so a small net value need not mean little
positive input or little extraction. This splitting applies after spatial
summation, rather than to positive and negative local power densities.

Both grids pass the predeclared $10^{-6}$ normalized balance tolerance,
and the new endfields agree with the stored reference endpoints within
the fixed $10^{-8}$ single-response-normalized tolerance. The reported
scaled two-grid differences of work contributions are below
$4.4\times10^{-4}$, within the prescribed $0.03$ consistency limit.
These are discrete numerical checks, not rigorous continuum error bounds;
the positive/negative power split also need not inherit a smooth fourth-order
convergence rate. The original nominal-mode boundary defect is unchanged.
Equal source-L2 budgets are still not equal work. Finally, the laboratory
source-work coefficient and $E_2$ belong to the frozen linear ansatz:
second-order background response, full nonlinear laboratory energy,
actuator losses and conversion efficiency are not determined. No complete
shutdown or long-time stability follows from this short-window balance.
\paragraph{Phase dependence of the second-pulse source work.}
For the same four source designs, we vary the constant phase $\theta$ of
the delayed pulse while keeping its source norm fixed. Real linearity of
the frozen evolution gives the quadratic source-work coefficient
\begin{equation}
 W_{E,\mathrm{src},2}(\theta)
 =c_0+c_1\cos\theta+s_1\sin\theta
       +c_2\cos(2\theta)+s_2\sin(2\theta).
 \label{eq:phase-work-five}
\end{equation}
Writing the measured quadrature coefficients as $(A,B,D,E,F)$, these are
$(c_0,c_1,s_1,c_2,s_2)=((D+E)/2,A,B,(D-E)/2,F/2)$.
The five coefficients are integrated from the first-pulse response and two
delayed quadrature responses, with a direct $45^\circ$ two-pulse control,
through $T=3$ after both source windows end ($\Delta+1<3$), using
$h=0.04,0.02$ and $\Delta t=h/8$. This is not a renewed $T=6$ energy history.
Thus the plotted curves are postprocessing of measured coefficients, not
thousands of independent evolution runs.

On both nominal grids, only B/T1 has a negative-work interval:
approximately $-66.53^\circ<\theta<57.03^\circ$.
For the fine grid its second-pulse net source work is
$-1.1224\times10^{-4}$ at $0^\circ$ and
$-3.9278\times10^{-5}$ in the direct $45^\circ$ control.
This broad nominal interval concerns signed net source work; it does not
assert equally small residual fields throughout the interval.
In particular, the previously reported residual field norms increase at
$45^\circ$ even though this work remains negative.
All four uniform \emph{constant-phase} means $c_0$ are positive;
for B/T1 the fine-grid mean is $1.1334\times10^{-4}$.
This phase ensemble is not a model of time-dependent phase noise.

Polynomial sign isolation treats the measured floating-point coefficients
as exact rational dyads. Its exactness applies to those nominal polynomials,
not to their discretization errors or to the continuum or nonlinear field
equations. The result concerns quadratic source work in the frozen linear
ansatz, with background exchange distinguished separately; it is neither a
full nonlinear laboratory-energy balance nor an actuator-efficiency claim.

\begin{figure}[htbp]
\centering
\includegraphics[width=\linewidth]{phase-work.pdf}
\caption{Second-pulse source work versus constant phase. Left: all four
designs; solid curves use $h=0.02$, dashed curves $h=0.04$.
Right: B/T1 with nominal fine-grid zero boundaries (dotted), its uniform
constant-phase mean (dash-dotted), and direct $45^\circ$ controls (symbols).
Work is displayed in units of $10^{-4}$. Shading denotes the nominal
negative-work window, not certified physical tolerances.}
\label{fig:phase-work}
\end{figure}
